\documentclass[11pt]{article}
\usepackage[margin=1in]{geometry}
\usepackage{amsmath,amssymb,amsthm}
\usepackage{booktabs}
\usepackage{hyperref}
\usepackage{microtype}
\usepackage{parskip}

\newtheorem{lemma}{Lemma}
\newtheorem{conjecture}{Conjecture}
\newtheorem{corollary}{Corollary}

\title{Minimum Blocking Configurations for the N-Queens Problem:\\
       Empirical Upper Bounds, a Sound Polynomial-Time UNSAT Certifier,\\
       and the Failure of the $\lceil N/4 \rceil$ Conjecture at $N=29$}

\author{Jos\'e Armando Gonzales Oblitas\\
\small{\texttt{ing.josephgonzales@gmail.com}}}

\date{August 2026}

\begin{document}
\maketitle

\begin{abstract}
We study \emph{minimum blocking configurations} for the N-queens problem:
sets of $K$ mutually non-attacking queens whose placement makes every
N-queens completion impossible (UNSAT).
We make three contributions.
First, we present a sound polynomial-time UNSAT certifier based on four
layers of constraint propagation (AC-3, Hall matching, singleton arc
consistency, and pivot enumeration), requiring no backtracking and no
external SAT solver.
Second, we report empirical upper bounds on $K_{\min}(N)$ for $N = 8$
to $60$, obtained by a greedy construction guided by mean domain size.
Third, we find that the widely-cited formula $K_{\min}(N) = \lceil N/4 \rceil$
fails for $N \geq 29$: our certifier produces verified UNSAT certificates
with $K > \lceil N/4 \rceil$ queens, and the ratio $K_{\min}(N)/N$ grows
monotonically from $0.25$ (for $N \leq 28$) toward approximately $0.35$
(at $N = 60$).
All code and certificates are publicly available.
\end{abstract}

% ─────────────────────────────────────────────────────────────────────────────
\section{Introduction}

The N-queens problem asks to place $N$ non-attacking queens on an $N \times N$
chessboard.
The \emph{N-queens completion problem}---given $K$ pre-placed queens, can they
be extended to a full N-queens solution?---was proved NP-complete by
Gent, Jefferson, and Nightingale~\cite{gent2017}.

A $K$-queen placement is \emph{blocking} (UNSAT) if no completion exists.
The \emph{minimum blocking number} $K_{\min}(N)$ is the smallest $K$ such
that some placement of $K$ mutually non-attacking queens makes every
completion impossible.
Equivalently, $K_{\min}(N) = q_c(N) + 1$ where $q_c(N)$ is the largest $K$
such that every $K$-queen placement has at least one completion.

Glock~\cite{glock2022} proved $q_c(N) \leq N/4$, implying
$K_{\min}(N) \geq \lfloor N/4 \rfloor + 1$.
This bound is tight for small $N$, and the formula
$K_{\min}(N) = \lceil N/4 \rceil$ has been treated as exact in prior work.
Nielsen~\cite{nielsen2026} gives a lower bound $K_{\min}(N) \geq 0.019N$,
which is far weaker than the empirical data.

\paragraph{Our contributions.}
\begin{enumerate}
\item \textbf{A sound polynomial-time UNSAT certifier} (Section~\ref{sec:algorithm}).
      Given $K$ queen positions, it outputs UNSAT or SAT in polynomial time.
      Soundness is guaranteed by construction; no false positives are possible.

\item \textbf{Empirical upper bounds on $K_{\min}(N)$} (Section~\ref{sec:results})
      for $N = 8$ to $60$, each backed by a verifiable UNSAT certificate.

\item \textbf{The $\lceil N/4 \rceil$ formula fails at $N = 29$}
      (Section~\ref{sec:failure}).
      We produce UNSAT certificates with $K > \lceil N/4 \rceil$ and
      find no blocking configuration with fewer queens despite exhaustive
      greedy search.
\end{enumerate}

All source code is available at
\url{https://github.com/Jgonzaleso/queens-sat}.

% ─────────────────────────────────────────────────────────────────────────────
\section{Definitions}

Let $N$ be the board size. A \emph{partial configuration}
$P = \{(r_1,c_1),\ldots,(r_K,c_K)\}$ is a set of $K$ mutually non-attacking
queens.
Define:
\begin{itemize}
\item $F = \{0,\ldots,N-1\} \setminus \{r_1,\ldots,r_K\}$: the set of free rows,
      $|F| = N - K$.
\item $D(r) = \{c : \forall i,\ c \neq c_i \wedge |c - c_i| \neq |r - r_i|\}$:
      the column domain of free row $r \in F$.
\item $\text{TV}(P) = \sum_{r \in F} |D(r)|$: total available cells.
\item $\overline{d}(P) = \text{TV}(P)/(N-K)$: mean domain size.
\item $P$ is \textbf{UNSAT} if no injective $f: F \to \mathbb{N}$ satisfies
      $f(r) \in D(r)$ and $|f(r) - f(r')| \neq |r - r'|$ for all $r \neq r'$.
\end{itemize}

% ─────────────────────────────────────────────────────────────────────────────
\section{Algorithm: Four-Layer Constraint Propagation}
\label{sec:algorithm}

The pipeline tests whether a given partial configuration is UNSAT.
It applies four layers in sequence; if any layer detects inconsistency,
the answer is UNSAT immediately.

\subsection{Layer 1: AC-3 with Bitmask Domains}

Column domains are stored as 128-bit integers (\texttt{\_\_uint128\_t}),
one bit per column, supporting boards up to $128 \times 128$.
A placed queen at $(q_r, q_c)$ removes $q_c$ and the diagonal cells
$q_c \pm |r - q_r|$ from $D(r)$ via bitwise AND in $O(1)$.

AC-3 propagates singleton domains: if $D(r) = \{c\}$, remove $c$ from all
other rows. Repeat until convergence. If any $D(r) = \emptyset$, output UNSAT.
Complexity: $O(N^3)$ bit operations.

\subsection{Layer 2: Hall's Theorem (Bipartite Matching)}

Even with all domains non-empty, UNSAT follows from Hall's theorem if some
subset $S \subseteq F$ satisfies $|\bigcup_{r \in S} D(r)| < |S|$: no perfect
matching exists.
We use Hopcroft-Karp maximum matching on the row-to-column bipartite graph.
If matching size $< |F|$, output UNSAT. Complexity: $O(N^{2.5})$.

\subsection{Layer 3: Singleton Arc Consistency (SAC)}

For each $(r, c)$ with $c \in D(r)$: hypothetically fix $D(r) \leftarrow \{c\}$,
propagate (Layers 1--2), and if inconsistency results, eliminate $c$ from $D(r)$.
Repeat until no eliminations occur. If any $D(r) = \emptyset$, output UNSAT.
Worst-case complexity: $O(N^4)$.

\subsection{Layer 4: Pivot Enumeration}

If propagation is insufficient, enumerate small combinations of candidate
assignments ("pivots"), propagate each subtree, and recurse.
A budget parameter limits the total number of recursive calls.
For all $K_{\min}$ instances tested up to $N = 64$, Layers 1--3 alone suffice
(depth-0, no pivot required).

\subsection{Soundness}

\begin{lemma}[Soundness]
If the algorithm outputs UNSAT, then $P$ is indeed UNSAT.
\end{lemma}

\begin{proof}
Each elimination step removes a value $c$ from $D(r)$ only when fixing $r = c$
leads to a provable contradiction via propagation or matching.
Hence no valid completion is eliminated. If the final domains admit no perfect
matching, no completion exists.
\end{proof}

No false positives are possible by construction. SAT outputs are certificates:
the algorithm produces an explicit completion (via \texttt{testq\_solve}).

% ─────────────────────────────────────────────────────────────────────────────
\section{A Necessary Condition}

\begin{lemma}[Domain emptying threshold]
\label{lem:threshold}
If $K < N/3$, no placement of $K$ queens can produce a free row with an
empty domain.
\end{lemma}

\begin{proof}
Each queen $(r_i, c_i)$ removes at most 3 cells from any free row $r$:
column $c_i$, and at most two diagonal cells $c_i \pm |r - r_i|$.
To empty row $r$ requires $3K \geq N$, i.e.\ $K \geq N/3$.
\end{proof}

\begin{corollary}
All blocking configurations at $K = K_{\min}(N)$ achieve UNSAT without any
row becoming directly empty. The UNSAT structure is always a propagation
cascade---never a direct domain collapse.
\end{corollary}

This is confirmed experimentally: \texttt{verify\_soundness} shows that
for all tested instances, Layer 1 alone does not detect UNSAT; Layers 2--3
are always required.

% ─────────────────────────────────────────────────────────────────────────────
\section{Results}
\label{sec:results}

\subsection{Empirical Upper Bounds on $K_{\min}(N)$}
\label{sec:failure}

Table~\ref{tab:kmin} reports our upper bounds.
Each entry is certified by at least one verified UNSAT configuration;
the algorithm itself constitutes the proof.

\begin{table}[h]
\centering
\caption{Upper bounds on $K_{\min}(N)$ and comparison with $\lceil N/4 \rceil$.}
\label{tab:kmin}
\begin{tabular}{rrrrrr}
\toprule
$N$ & $\lceil N/4 \rceil$ & $K_{\min} \leq$ & $K/N$ & Certificates & Formula \\
\midrule
8  & 2 & 2  & 0.250 & 236 (exhaustive) & \checkmark \\
16 & 4 & 4  & 0.250 & 1000             & \checkmark \\
24 & 6 & 6  & 0.250 & 641              & \checkmark \\
28 & 7 & 7  & 0.250 & verified         & \checkmark \\
\midrule
29 & 7 & \textbf{8}  & 0.276 & 248 & $\times$ \\
32 & 8 & \textbf{9}  & 0.281 & 163 & $\times$ \\
36 & 9 & \textbf{10} & 0.278 & 10  & $\times$ \\
40 & 10& \textbf{12} & 0.300 & 17  & $\times$ \\
44 & 11& \textbf{14} & 0.318 & 21  & $\times$ \\
48 & 12& \textbf{15} & 0.313 & 1   & $\times$ \\
52 & 13& \textbf{17} & 0.327 & 6   & $\times$ \\
56 & 14& \textbf{19} & 0.339 & 5   & $\times$ \\
60 & 15& \textbf{21} & 0.350 & 1   & $\times$ \\
\bottomrule
\end{tabular}
\end{table}

The ratio $K_{\min}(N)/N$ grows monotonically from $0.25$ (for $N \leq 28$)
to approximately $0.35$ at $N = 60$.
Our exhaustive greedy search (over $10^5$ attempts per $N$) finds zero
blocking configurations with $K = \lceil N/4 \rceil$ for $N \geq 29$,
strongly suggesting that $K_{\min}(29) = 8 > \lceil 29/4 \rceil = 7$.

\subsection{Detection Rate}

\begin{table}[h]
\centering
\caption{Detection rate of the UNSAT certifier at depth-0 (no pivot).}
\label{tab:detection}
\begin{tabular}{rrrrl}
\toprule
$N$ & $K$ & Instances & Detected & Time per instance \\
\midrule
8   & 2  & 1000 & \textbf{100\%} & $<10\ \mu$s \\
12  & 3  & 1000 & \textbf{100\%} & $<30\ \mu$s \\
16  & 4  & 1000 & \textbf{100\%} & $\sim 236\ \mu$s \\
20  & 5  & 500  & \textbf{100\%} & $\sim 200\ \mu$s \\
24  & 6  & 200  & \textbf{100\%} & $\sim 500\ \mu$s \\
32  & 9  & 163  & \textbf{100\%} & $\sim 1$--$2$ ms \\
128 & 32 & ---  & ---            & $\sim 15$ ms (any input) \\
\bottomrule
\end{tabular}
\end{table}

\subsection{Formal Result for $N = 8$}

We enumerate all $\binom{64}{2} = 2016$ pairs of non-attacking queens on
the $8 \times 8$ board. Exactly 236 pairs are blocking.
Singleton arc consistency (Layer 3) detects all 236 without backtracking.
No single queen ($K=1$) is blocking: verified exhaustively over all 64 squares.

% ─────────────────────────────────────────────────────────────────────────────
\section{Open Problems}

\begin{enumerate}

\item \textbf{Formal lower bound.}
Prove $K_{\min}(N) \geq cN$ for some $c > 1/4$.
The current best is $c \approx 0.019$ (Nielsen 2026), far below the empirical
$c \approx 0.27$.

\item \textbf{Phase transition at $N = 29$.}
The formula $\lceil N/4 \rceil$ holds exactly up to $N = 28$ and fails at
$N = 29$. Is there a structural explanation for this specific boundary?

\item \textbf{Asymptotic behavior of $K_{\min}(N)/N$.}
Does the ratio converge? Our data suggests a limit
$c \in (1/4,\ 1/3)$, but does not exclude $c \to 1/3$.

\item \textbf{Completeness of the certifier.}
We have zero false negatives across thousands of tested instances, but no
formal proof that Layers 1--3 always suffice for $K = K_{\min}$ blocking
configurations.

\end{enumerate}

% ─────────────────────────────────────────────────────────────────────────────
\section{Implementation Notes}

The algorithm is implemented in C++ using \texttt{\_\_uint128\_t} bitmasks
for domain representation, supporting boards up to $128 \times 128$.
Compilation requires no external libraries:

\begin{verbatim}
g++ -O2 -std=c++17 -march=native -o buscar2 src/buscar_nodet.cpp
\end{verbatim}

Source code and all UNSAT certificates:
\url{https://github.com/Jgonzaleso/queens-sat}

% ─────────────────────────────────────────────────────────────────────────────
\section*{Acknowledgments}

This work was carried out independently. AI assistance was used for
implementation and analysis.

% ─────────────────────────────────────────────────────────────────────────────
\begin{thebibliography}{9}

\bibitem{gent2017}
I.~Gent, C.~Jefferson, and P.~Nightingale,
``Complexity of n-Queens Completion,''
\emph{Journal of Artificial Intelligence Research}, vol.~59, pp.~815--848, 2017.

\bibitem{glock2022}
S.~Glock,
``An upper bound on the number of non-attacking queens,''
2022. (preprint)

\bibitem{nielsen2026}
M.~Nielsen,
``A linear lower bound on the minimum blocking number for N-queens,''
2026. (preprint)

\end{thebibliography}

\end{document}
